% Section 7.8, translated from sv/chapters/07/exercises.tex.

\section{Exercises}
\label{c7:sec:uppgifter}

\begin{aside}
\textbf{How to work with the exercises.} The vector exercises in Part~1 are done during the lesson:
first on your own, a few minutes per exercise, then together. Part~2 is extra exercises for
further practice after the lesson. Most of them can be done in your head or with pencil and paper.

\facitnotis{L07}
\end{aside}

% ------------------------------------------------------------------------------------------
\uppgiftsdel{Part 1}{Vector exercises}

\uppgift{1.1}{Addition and scalar multiplication}
The vectors $\mathbf{u} = (3, 2)$ and $\mathbf{v} = (4, 6)$ are given. Find the coordinates
of:
\begin{delar}(4)
  \task $\mathbf{u} + \mathbf{v}$
  \task $2\mathbf{u} + 3\mathbf{v}$
  \task $\mathbf{u} - \mathbf{v}$
  \task $-2\mathbf{u} - 4\mathbf{v}$
\end{delar}

\uppgift{1.2}{Calculating with three vectors}
The vectors $\mathbf{a} = (1, 4)$, $\mathbf{b} = (-4, 2)$ and $\mathbf{c} = (3, -4)$ are
given. Find the coordinates of:
\begin{delar}(4)
  \task $\mathbf{a} + \mathbf{b} + \mathbf{c}$
  \task $\mathbf{a} - \mathbf{b} - \mathbf{c}$
  \task $2\mathbf{a} + 4\mathbf{b}$
  \task $2\mathbf{a} - 4\mathbf{b}$
\end{delar}

\uppgift{1.3}{Magnitude}
The vectors $\mathbf{a} = (2, 1)$ and $\mathbf{b} = (1, -2)$ are given. Find, in exact form,
the length of:
\begin{delar}(4)
  \task $\mathbf{a}$
  \task $3\mathbf{a} + \mathbf{b}$
  \task $\mathbf{a} - \mathbf{b}$
  \task $3\mathbf{a} - 4\mathbf{b}$
\end{delar}

\uppgift{1.4}{Solve the equations}
Find $x$ and $y$.
\begin{delar}(2)
  \task $(x, y) = 5(3, 1) + 2(4, -5)$
  \task $(15, 12) = (x, 10) + (23, y)$
\end{delar}

\uppgift{1.5}{Parallel vectors}
If $\mathbf{u} = k \cdot \mathbf{v}$, the vectors $\mathbf{u}$ and $\mathbf{v}$ are parallel.
Which of the following vectors are parallel to $\mathbf{u} = (2, -3)$?
\begin{delar}(4)
  \task $(6, -9)$
  \task $(-8, 16)$
  \task $(10, 15)$
  \task $(-12, 18)$
\end{delar}

\uppgift{1.6}{Finding unknown coefficients}
The vectors $\mathbf{u} = (-2, 3)$ and $\mathbf{v} = (3, -1)$ are given. Find the numbers $a$
and $b$ such that
\[
  a\mathbf{u} + \mathbf{v} = (0, b).
\]

% ------------------------------------------------------------------------------------------
\uppgiftsdel{Part 2}{Extra exercises}
The exercises below are extra practice, best done on your own after the lesson.

\uppgift{2.1}{Distance between points}
Calculate the distance between the points:
\begin{delar}(3)
  \task $(0, 0)$ and $(3, 4)$
  \task $(1, 2)$ and $(4, 6)$
  \task $(-2, 1)$ and $(2, -2)$
  \task $(5, -1)$ and $(5, 7)$
  \task $(-3, 2)$ and $(1, 2)$
\end{delar}

\uppgift{2.2}{The vector between two points}
The vector from the point $A$ to the point $B$ is given by
$\overrightarrow{AB} = (x_B - x_A, y_B - y_A)$.
\begin{parts}
  \item $A = (1, 2)$ and $B = (5, 5)$. Find $\overrightarrow{AB}$ and
    $\abs{\overrightarrow{AB}}$.
  \item $A = (-1, 4)$ and $B = (2, 0)$. Find $\overrightarrow{AB}$ and
    $\abs{\overrightarrow{AB}}$.
  \item Find $\overrightarrow{BA}$ in a) and compare it with $\overrightarrow{AB}$.
  \item What holds in general between $\overrightarrow{AB}$ and $\overrightarrow{BA}$?
\end{parts}

\uppgift{2.3}{Magnitude}
Find $\abs{\mathbf{u}}$ in exact form and as an approximate value to two decimal places.
\begin{delar}(3)
  \task $\mathbf{u} = (6, 8)$
  \task $\mathbf{u} = (-5, 12)$
  \task $\mathbf{u} = (1, 1)$
  \task $\mathbf{u} = (-3, -3)$
  \task $\mathbf{u} = (0, -7)$
\end{delar}

\uppgift{2.4}{Angle with the $x$-axis}
Find the angle $v$ with the positive $x$-axis. State which quadrant the vector lies in and whether
a quadrant correction is needed.
\begin{delar}(5)
  \task $(1, 1)$
  \task $(-1, 1)$
  \task $(-1, -1)$
  \task $(1, -1)$
  \task $(4, 3)$
\end{delar}

\uppgift{2.5}{From magnitude and angle to components}
A vector of length $\abs{\mathbf{u}}$ at the angle $v$ to the positive $x$-axis has the components
\[
  u_x = \abs{\mathbf{u}} \cos v, \qquad u_y = \abs{\mathbf{u}} \sin v.
\]
Find the components:
\begin{delar}(2)
  \task $\abs{\mathbf{u}} = 10$, $v = 30^{\circ}$
  \task $\abs{\mathbf{u}} = 5$, $v = 90^{\circ}$
  \task $\abs{\mathbf{u}} = 4$, $v = 180^{\circ}$
  \task $\abs{\mathbf{u}} = 2$, $v = 225^{\circ}$
  \task* Check in a) that $\sqrt{u_x^2 + u_y^2} = 10$.
\end{delar}

\uppgift{2.6}{Scalar product}
Calculate $\mathbf{u} \cdot \mathbf{v}$.
\begin{delar}(2)
  \task $\mathbf{u} = (2, 3)$, $\mathbf{v} = (4, 1)$
  \task $\mathbf{u} = (1, -2)$, $\mathbf{v} = (3, 5)$
  \task $\mathbf{u} = (2, 4)$, $\mathbf{v} = (-6, 3)$
  \task $\mathbf{u} = (5, 0)$, $\mathbf{v} = (0, 7)$
  \task* What do c) and d) have in common, and what does it say about the vectors?
\end{delar}

\uppgift{2.7}{Perpendicular vectors}
Two vectors are perpendicular (orthogonal) if $\mathbf{u} \cdot \mathbf{v} = 0$.
\begin{parts}
  \item Are $(3, 4)$ and $(-4, 3)$ perpendicular?
  \item Are $(1, 2)$ and $(2, 1)$ perpendicular?
  \item Find the number $k$ such that $(2, k)$ is perpendicular to $(6, -3)$.
  \item Give a vector that is perpendicular to $(5, -2)$.
\end{parts}

\uppgift{2.8}{The angle between two vectors}
Calculate the angle $\alpha$ between the vectors using
$\cos\alpha = \dfrac{\mathbf{u} \cdot \mathbf{v}}{\abs{\mathbf{u}}\abs{\mathbf{v}}}$.
\begin{delar}(2)
  \task $\mathbf{u} = (1, 0)$, $\mathbf{v} = (1, 1)$
  \task $\mathbf{u} = (3, 0)$, $\mathbf{v} = (0, 4)$
  \task $\mathbf{u} = (2, 1)$, $\mathbf{v} = (-1, 3)$
  \task $\mathbf{u} = (2, 2)$, $\mathbf{v} = (-3, -3)$
\end{delar}

\uppgift{2.9}{Unit vectors}
A unit vector has length $1$ and is obtained as
$\hat{\mathbf{u}} = \dfrac{\mathbf{u}}{\abs{\mathbf{u}}}$.
\begin{parts}
  \item Find the unit vector in the direction $(3, 4)$.
  \item Check that your answer to a) really has length $1$.
  \item Find the unit vector in the direction $(-6, 8)$.
  \item Find a vector of length $15$ in the same direction as $(3, 4)$.
  \item Find a vector of length $15$ in the opposite direction to $(3, 4)$.
\end{parts}

\uppgift{2.10}{Resultant}
Three forces act at the same point: $\mathbf{F}_1 = (4, 2)$, $\mathbf{F}_2 = (-1, 5)$ and
$\mathbf{F}_3 = (-3, -7)$.
\begin{parts}
  \item Calculate the resultant $\mathbf{F} = \mathbf{F}_1 + \mathbf{F}_2 + \mathbf{F}_3$.
  \item What does the result mean?
  \item Calculate the resultant and its magnitude if $\mathbf{F}_3$ is removed.
  \item Which vector must be added to $\mathbf{F}_1 + \mathbf{F}_2$ for the resultant to be zero?
\end{parts}

\uppgift{2.11}{Parallel and opposite vectors}
The vector $\mathbf{u} = (4, -6)$ is given.
\begin{parts}
  \item Which of $(2, -3)$, $(-8, 12)$, $(6, -4)$ and $(-2, 3)$ are parallel to
    $\mathbf{u}$?
  \item Which of them point in the opposite direction to $\mathbf{u}$?
  \item Find $t$ such that $(t, 9)$ is parallel to $\mathbf{u}$.
  \item Can a vector be parallel to itself? Justify your answer.
\end{parts}

\uppgift{2.12}{Voltages as phasors}
In an AC network, voltages are represented by vectors known as \textbf{phasors}. The voltages
across a resistor and an inductor in series are $\mathbf{U}_R = (12, 0)\,\text{V}$ and
$\mathbf{U}_L = (0, 5)\,\text{V}$.
\begin{parts}
  \item Calculate the total voltage $\mathbf{U} = \mathbf{U}_R + \mathbf{U}_L$.
  \item Calculate the magnitude $\abs{\mathbf{U}}$.
  \item Calculate the angle with the positive $x$-axis.
  \item Why is $\abs{\mathbf{U}}$ not equal to $12 + 5 = 17\,\text{V}$?
\end{parts}

\uppgift{2.13}{Rules for the scalar product}
The vectors $\mathbf{u} = (2, -1)$, $\mathbf{v} = (3, 4)$ and $\mathbf{w} = (-1, 2)$ are
given.
\begin{parts}
  \item Calculate $\mathbf{u} \cdot \mathbf{v}$ and $\mathbf{v} \cdot \mathbf{u}$. What holds?
  \item Calculate $\mathbf{u} \cdot (\mathbf{v} + \mathbf{w})$ and
    $\mathbf{u} \cdot \mathbf{v} + \mathbf{u} \cdot \mathbf{w}$. What holds?
  \item Calculate $\mathbf{u} \cdot \mathbf{u}$ and compare it with $\abs{\mathbf{u}}^2$.
  \item What is $\mathbf{u} \cdot \mathbf{u}$ in terms of $\abs{\mathbf{u}}$, in general?
\end{parts}

\uppgift{2.14}{Find the error}
Each line contains a common mistake with vectors. Explain the mistake and give the correct answer.
\begin{parts}
  \item $\abs{(3, 4)} = 3 + 4 = 7$
  \item $2(3, -1) = (6, -1)$
  \item $(2, 3) \cdot (4, 5) = (8, 15)$
  \item The vector $(-3, 4)$ makes the angle
    $\arctan\!\left(\dfrac{4}{-3}\right) \approx -53.1^{\circ}$ with the positive $x$-axis.
  \item $\abs{-2\mathbf{u}} = -2\abs{\mathbf{u}}$
\end{parts}
